\section{\Egglog}
\label{sec:overview}

\Egglog is a logic programming language 
 that bears many similarities to Datalog,
 and it also incorporates features 
 that allow for program optimization and verification as in equality saturation.
In this section, 
 we approach \egglog by example,
 starting from the Datalog perspective
 and adding features until 
 it subsumes equality saturation.

\subsection{Datalog in \Egglog}

\begin{figure}
    \centering
    \hfill
    \begin{subfigure}[t]{0.40\linewidth}
        \begin{lstlisting}[language=egglog, numbers=left,escapechar=!,xleftmargin=1em]
(relation edge (i64 i64)) !\label{lst:edge-decl}!
(relation path (i64 i64)) !\label{lst:path-decl}!
            
(rule ((edge x y)) !\label{lst:tc-start}!
      ((path x y)))
(rule ((path x y) (edge y z))
      ((path x z))) !\label{lst:tc-end}!

(edge 1 2) !\label{lst:edge-1}!
(edge 2 3) !\label{lst:edge-2}!
(edge 3 4) !\label{lst:edge-3}!

(run) 
(check (path 1 4)) ;; succeeds
\end{lstlisting}
        \caption{Reachability in the classic Datalog style.}
        \label{fig:egglog-is-datalog:a}
    \end{subfigure}
    \hfill
    \begin{subfigure}[t]{0.55\linewidth}
        \begin{lstlisting}[language=egglog, numbers=left]
(function edge (i64 i64) i64)
(function path (i64 i64) i64 :merge (min old new))
            
(rule ((= (edge x y) len))
      ((set (path x y) len)))
(rule ((= (path x y) xy) (= (edge y z) yz))
      ((set (path x z) (+ xy yz))))

(set (edge 1 2) 10)
(set (edge 2 3) 10)
(set (edge 1 3) 30)

(run) 
(check (path 1 3)) ;; prints "20"
\end{lstlisting}
        \caption{Reachability including shortest path length.}
        \label{fig:egglog-is-datalog:b}
    \end{subfigure}
    \caption{
        \Egglog supports classic Datalog programs like reachability
         written in the natural way.
        Functions and \lstinline{:merge}
         allow \egglog to support Datalog with lattices similar to tools
         like Flix~\cite{flix} or Ascent~\cite{ascent}.
    }
    \label{fig:egglog-is-datalog}
\end{figure}

\Egglog uses a concrete syntax based on s-expressions,
 but despite this surface-level difference,
 readers familiar with Datalog should find many 
 \egglog programs familiar.
The program in \autoref{fig:egglog-is-datalog:a}
 computes the transitive closure of a graph,
 just like the Datalog program in \autoref{fig:datalog-tc-code}.
It first declares two relations of pairs of 64-bit integers 
 (\lstinline{i64} is one of \egglog's base types).
The \lstinline{edge} relation stores the edges of a graph 
 and is initially populated manually on lines~\ref{lst:edge-1}-\ref{lst:edge-3}.
The \lstinline{path} relation is populated by the rules on lines~\ref{lst:tc-start}-\ref{lst:tc-end}.
These rules compute the transitive closure of the \lstinline{edge} relation.
Finally, the last two lines execute the program and check that 
 there is a path from 1 to 4.
% \autoref{fig:egglog-is-datalog:a} 
%  declares two relations, 
%  declares two rules,
%  populates the \lstinline{edge} relation with three facts,
%  and then runs the program.%
% \footnote{
%   In Datalog, running the program is usually implicit,
%   but \egglog requires an explicit \lstinline{(run)} command.
%   This becomes more useful later when 
%   we may not desire to run the program all the way to fixpoint.
% }

Let us take a closer look at the second rule in \autoref{fig:egglog-is-datalog:a}:
 it states that if there is a path from $x$ to $y$ 
 and an edge from $y$ to $z$, 
 then there is a path from $x$ to $z$.
In \egglog, a rule has two parts: 
 a \emph{query} and a list of \emph{actions}.%
\footnote{
Note that this is backwards from the
 more traditional Datalog syntax:
 \lstinline{path(X, Z) :- path(X, Y), edge(Y, Z).}
An \egglog rule's query and actions are analogous 
 to the body and head of a Datalog rule, respectively.
}
The query is a set of \emph{patterns},
 all of which must match for the rule to fire.
If all patterns do match, the query binds each variable to a value.
The actions dictate what happens when the rule fires,
 and they can use the variables that are bound by the query.
Typically, as in this example, the actions assert
 new facts to be added to the database.

% \yz{talk about \lstinline[language=egglog]{define}}

\subsection{Functions and \lstinline{:merge}}
\label{sec:merge}

Unlike traditional Datalog, 
 \egglog stores data as partial functions rather than relations.
A relation in \egglog actually 
 desugars to a function whose return type is the built-in unit type.
To model a unary relation $R$,
 we can use a function $f_R$ to unit such that
 $f_R(x) = () \text{ if } x \in R \text{ else undefined}$.
While a Datalog program's rules add tuples to relations,
 \egglog's functions become defined for more and more tuples
 over the course of a program's execution,
 a concept that we will explore in more detail in \autoref{sec:semantics}.
Every user-defined function in \egglog is
 backed by a map (as opposed to a set in in Datalog).
Crucially, the map enforces the functional dependency from inputs to outputs.
In other words, a function maps each input to a unique output.
Throughout this paper, 
 we will use the term ``table'' to refer to either
 the backing map of an \egglog function
 or the backing set of a Datalog relation.

Consider the program in \autoref{fig:egglog-is-datalog:b}
 that computes the length of shortest path between all nodes.
Ignoring the \lstinline{:merge} declaration for now,
 the program is substantially similar 
 to the reachability program in \autoref{fig:egglog-is-datalog:a},
 but it uses functions instead of relations.
The program defines 
 \lstinline{edge} and \lstinline{path} functions to \lstinline{i64}
 rather than functions to unit (i.e., relations).
The first rule (the base case) in \autoref{fig:egglog-is-datalog:b}
 is similar to before: 
 it says that an edge of length $\mathsf{len}$ from $x$ to $y$
 implies there is a path of (at most) length $\mathsf{len}$ from $x$ to $y$.
The query uses
 \lstinline{=} to bind the output
 of the \lstinline{edge} function to the variable \lstinline{len}.
% A query is a list of unit-typed facts,
%  and the \lstinline{=} operator 
%  is a partial function from pairs of some type to unit.
In the action, 
 we see the \lstinline{set} construct,
 which asserts
 that a function maps some arguments to a given value.
The action in the 
 analogous rule in \autoref{fig:egglog-is-datalog:a}
 desugars to \lstinline{(set (path x y) ())}.
Note that if the arguments are already mapped to a value, 
 we need to reconcile the old value with the new one
 to preserve the functional dependency.
This is resolved by the \lstinline{:merge} declaration which 
 we will describe next.

The second rule in \autoref{fig:egglog-is-datalog:b}
 is the transitive case,
 and here we see the purpose of the \lstinline{:merge} declaration.
This rule says that if there is a path from $x$ to $y$ of length $xy$,
 and an edge from $y$ to $z$ of length $yz$,
 then there is a path from $x$ to $z$ of length $xy + yz$.
But what if the function \lstinline{path} is already defined
 on the arguments $x$ and $z$?
Functions must map equivalent arguments to unique output,
 so the \lstinline{:merge} declaration tells \egglog 
 how to resolve this conflict.
Given the facts later in \autoref{fig:egglog-is-datalog:b},
 the program will discover two paths from $1$ to $3$:
 the single edge with length $30$ 
 will be discovered first,
 and then two-edge path with length $20$.
When \lstinline{(set (path 1 3) (+ 10 10))} 
 is executed,
 \egglog must come up with a single value to map \lstinline{(path 1 3)} to.
To do this,
 it evaluates the expression given after \lstinline{:merge} 
 in the function's declaration
 with \lstinline{old} and \lstinline{new} bound to the old and new values,
 respectively.
% The default \lstinline{:merge} 
%  for functions that output base types 
%  crashes the program if a function called with the same arguments 
%  is set to two different values.
In this case,
 \lstinline{path}'s \lstinline{:merge} expression
 simply takes the minimum of the two path lengths.
It can be viewed as the join operator,
 which takes the supremum of a set of values, 
 of the min lattice over \lstinline{i64} 
 where the partial order is $x\sqsubseteq y\iff x\geq y$.
This is similar to the lattice semantics of Flix~\citep{flix},
 which also enforces functional dependency by taking the join
 over the old and new values in some lattice.
% While this is itself useful for many Datalog programs,
%  \autoref{fig:egglog-unification:b} shows how the lattice semantics corresponds to
%  and actually generalize \eclass analyses in equality saturation.
% \mw{Does it? Jrw is also confused about this.}
% \of{Is there an opportunity to say how egglog does better than just datalog because it only keeps the minimum path?}
However, \egglog does not restrict the \lstinline{:merge} expression
 to only join operations over lattices.
In the following sections we will show how a \lstinline{:merge} expression
 that unifies values naturally gives rise to equality saturation.

% \mw{something about lattices and flix here}
% \yz{
%     We should point out the beautiful connection between 
%     this lattice semantics of Datalog and e-class analysis
%     in the overview, or in the background section.
% }

\subsection{Sorts and Equality}

\Egglog gives the user the ability to declare new
 \emph{uninterpreted} sorts,
 and functions use these new sorts as inputs or outputs.
Crucially, values of user-defined sorts (as opposed to base types)
 can be \emph{unified} by the \lstinline{union} action.
\lstinline{union}-ing two values makes them point to the same element 
 in the underlying universe of uninterpreted sorts.
In other words, 
 values that have been unified are essentially indistinguishable
 to \egglog, 
 and all unified variables can be substituted for the same pattern variable.

\begin{figure}
    \centering
    \begin{subfigure}[t]{0.40\linewidth}
        \begin{lstlisting}[language=egglog]
(sort Node)
(function mk (i64) Node)
(relation edge (Node Node))
(relation path (Node Node))
            
(rule ((edge x y))
      ((path x y)))
(rule ((path x y) (edge y z))
      ((path x z)))

(edge (mk 1) (mk 2))
(edge (mk 2) (mk 3))
(edge (mk 5) (mk 6))

(union (mk 3) (mk 5))
(run)
(check (edge (mk 3) (mk 6)))
(check (path (mk 1) (mk 6)))
\end{lstlisting}
        \caption{Combining nodes with unification}
        \label{fig:egglog-unification:a}
    \end{subfigure}
    \begin{subfigure}[t]{0.55\linewidth}
        \begin{lstlisting}[language=egglog]
(datatype Math
  (Num i64)
  (Var String)
  (Add Math Math)
  (Mul Math Math))

;; expr1 = 2 * (x + 3)
(define expr1 (Mul (Num 2) (Add (Var "x") (Num 3)))) 
;; expr2 = 6 + 2 * x
(define expr2 (Add (Num 6) (Mul (Num 2) (Var "x"))))

(rewrite (Add a b)         (Add b a))
(rewrite (Mul a (Add b c)) (Add (Mul a b) (Mul a c)))
(rewrite (Add (Num a) (Num b)) (Num (+ a b)))
(rewrite (Mul (Num a) (Num b)) (Num (* a b)))

(run)
(check (= expr1 expr2))
\end{lstlisting}
        \caption{Basic equality saturation}
        \label{fig:egglog-unification:b}
    \end{subfigure}
    \caption{
        Unification and EqSat in \egglog.
    }
    \label{fig:egglog-unification}
\end{figure}

Consider an enhanced version of path reachability 
 in \autoref{fig:egglog-unification:a},
 where we use unification to implement node contraction
 (sometimes called vertex contraction).
The program declares a new sort \lstinline{Node},
 which is necessary because base types (like \lstinline{i64})
 cannot be unified.
The \lstinline{mk} function is the sole constructor of \lstinline{Node}s.
After the rule declarations (same as in \autoref{fig:egglog-is-datalog:a})
 and some edge assertions,
 we see our first \lstinline{union} action,
 which takes two arguments
 of the same user-defined sort and unifies them.
Now that nodes $3$ and $5$ are unified,
 running the rules will indeed find a path from $1$ to $6$, 
 a path that did not exist before the unification.

In \egglog, users define uninterpreted sorts.
A sort is a set of opaque integer values
 called \emph{ids} and an equivalence relation over those ids.
The equivalence relation is implemented
 with a union-find data structure~\cite{unionfind}
 that can \emph{canonicalize} ids;
 two ids are equivalent iff 
 they canonicalize to the same id.
Equivalent ids are considered indistinguishable by \egglog.
In fact,
 \egglog ensures that all ids appearing in the database are canonical.
These ids corresponds to \eclass ids from the \eqsat perspective.

% \of{the following paragraph is a bit confusing, and comes out of nowhere. What's the best place to put this information about defaults and total functions?}
The second line of \autoref{fig:egglog-unification:a}
 declares \lstinline{mk}, 
 a function from \lstinline{i64} to \lstinline{Node}.
This looks like a constructor, for \lstinline{Node}s, 
 but it is just like any other function from \egglog's perspective;
 the \lstinline{mk} function is backed 
 by a map from \lstinline{i64}s to \lstinline{Node} ids.
In this program,
 we never query over the \lstinline{mk} function,
 but we do call it, treating it like a total function.
What is the value of \lstinline{(mk 1)},
 especially since we did not \lstinline{set} it to anything
 prior to calling it?
Functions in \egglog can be imbued with a 
 \lstinline{:default} expression that extends
 the partial function as defined by the underlying map to be total.
Calling a function \lstinline{(f x)} will first see if
 the map for function \lstinline{f} defines 
 an output for \lstinline{x}.
If so, it returns that output.
Otherwise, 
 \egglog evaluates the \lstinline{:default} expression,
 stores the result in the map,
 and returns it.
Unless otherwise specified,
 the \lstinline{:default} for functions that output a
 user-defined sort is to create an equivalence class in the union-find
 and return its id
 (the ``make-set'' operation);
 for base types the default \lstinline{:default} 
 is to crash the program.
In other words,
 calling a function that outputs a user-defined sort
 is essentially a ``get or make-set'' operation.
% \yz{For sorts, :merge being union and :default being union is more like a must, 
%  which is implied by the semantics in \autoref{sec:semantics}.
%  This seem to imply the default can be anything}

The upcoming subsection will discuss how these features enable
 equality saturation,
 but \autoref{sec:points-to} will demonstrate how
 the canonicalizing union-find is useful even in a domain
 where Datalog is traditionally strong: pointer analysis.

\subsection{Terms and Equality Saturation}\label{sec:eqsat}

% \yz{Maybe good to come up with a throughout example that is used both here and in the introduction.}

In \autoref{fig:egglog-unification:a},
 the \lstinline{Node} sort only has a single constructor, 
 \lstinline{mk}, which takes an \lstinline{i64}.
\egglog also supports functions that take user-defined sorts as inputs.
In this way, terms are easily constructed in \egglog.
Combined with the built-in equivalence relation,
 this term representation directly supports equality saturation in \egglog.

Consider \autoref{fig:egglog-unification:b},
 where we define a datatype \lstinline{Math} 
 that represents a simple language of arithmetic expressions.
The \lstinline{datatype} construct is sugar for
 a \lstinline{sort} declaration
 and a \lstinline{function} declaration for each constructor.
Each constructor is a function that returns a value of type \lstinline{Math},
 and its \lstinline{:default} behavior creates a fresh id
 as described above
 (we will get to its \lstinline{:merge} behavior shortly).
Now we can create terms by just nesting function calls.
The \lstinline{define} statements do just that, 
 creating two terms that we will later prove are equivalent.
These statements actually create nullary (constant) functions;
 \lstinline{(define x e)} desugars to
 \lstinline{(function x () T) (set (x) e)} 
 where \lstinline{T} is the type of \lstinline{e}.
Evaluating these terms
 adds them to the database (if not already present)
 thanks to the \lstinline{:default} behavior of the constructors.

Term rewriting in equality saturation
 has two important qualities:
 (1) pattern matching is done \emph{modulo equality} and
 (2) rewriting is \emph{non-destructive}, i.e., 
     it only adds information to the e-graph/database.
\Egglog meets both of these criteria:
 (1) all queries are performed modulo equality since
     \egglog canonicalizes the database with respect 
     to its built-in equivalence relation,
     and
 (2) \egglog rules (like standard Datalog rules) only add information
     to the database.
\Egglog provides the 
 \lstinline{rewrite} statement
 to simplify creating
 equality saturation rewrite rules.
A \lstinline{(rewrite p1 p2)} statement
 desugars to a \lstinline{rule} that queries
 for \lstinline{p1}, binds it to some variable,
 and \lstinline{union}s the variable with \lstinline{p2}:
 \lstinline{(rule ((= __var p1)) ((union __var p2)))}.

The program in \autoref{fig:egglog-unification:b}
 proves \lstinline{expr1} equivalent to \lstinline{expr2}
 using two uninterpreted rewrites
 and two that interpret the \lstinline{Add} and \lstinline{Mul} functions
 using the built-in \lstinline{+} and \lstinline{*} functions over \lstinline{i64}.
\eqsat frameworks like \egg 
 require the user to separate the uninterpreted rules
 from the interpreted part from the computed part using
 an \eclass analysis~\cite{egg}.
\egglog uses rules for both.

Like other functions that output user-defined sorts,
 the \lstinline{Math} constructors' \lstinline{:merge} behavior
 is to union the two ids.
Combined with \egglog's canonicalization,
 this means that the built-in equivalence relation is also
 a congruence relation with respect to these functions.
Consider the following map for the
 \lstinline{Add} function:
 $\{ (a, b) \mapsto c, (a, d) \mapsto e \}$.
If we \lstinline{union} $b$ and $d$ 
 such that $b$ is now canonical,
 canonicalizing the database reveals a violation 
 of the functional dependency 
 from \lstinline{Add}'s inputs to its output:
 $(a, b) \mapsto c$ but also $(a, b) \mapsto e$.
To resolve the conflict,
 \egglog invokes the \lstinline{:merge} expression
 of the \lstinline{Add} function,
 which in this case \lstinline{union}s $c$ and $e$.

After running the rules,
 the final line \lstinline{check}s 
 that \lstinline{expr1} and \lstinline{expr2} are now equivalent.
The type of both \lstinline{expr1} and \lstinline{expr2} 
 is \lstinline{Math}---a user-defined sort---%
 so the underlying value of the expressions are both ids.
Since \egglog canonicalizes the database,
 the \lstinline{check} is implemented 
 with simple equality on the ids.
\egglog supports optimization as well as verification;
 the \lstinline{extract} command prints the 
 smallest term equivalent to its given input. 

% \yz{TODO: add e-class analyses to this EqSat overview}

\subsection{Beyond \eqsat}

\Egglog is not limited to just Datalog or \eqsat;
 the combination allows for possibilities outside the reach of either tool.
The combining nodes example 
 from \autoref{fig:egglog-is-datalog:a}
 hints at the power of unification in Datalog,
 and \autoref{sec:points-to} takes this 
 further by implementing a unification-based pointer analysis in \egglog.
In \autoref{sec:herbie},
 we go the other way, 
 implementing several Datalog-like anaylses
 to assist an \eqsat-powered term rewriting system.

But \egglog goes beyond these applications;
 we
 describes more \egglog pearls in the full version \citep{egglog-preprint},
 including functional programming,
 type analyses for the simply typed lambda calculus in equality saturation,
 type inference for Hindley-Milner type systems,
 multivariable equational solving,
 and matrix algebra optimization with Kronecker products.
% While Datalog and \eqsat engines may struggle to express these applications,
%  users can express them straightforwardly in \egglog
These pearls hint at the potential novel applications in program optimizations and analyses using \egglog in the future.
% While writing these applications in \egglog is straightforward,
%  writing them in Datalog or \egglog can be tricky if not impossible.

% But \egglog does not stop here.
% \autoref{sec:pearls} is dedicated to demonstrating this point.
% \autoref{sec:funprog} shows how to express 
%  first-order functional programs in \egglog
%  in a natural way.
% While top-down semantic reasoning is difficult to express in \eclass analyses,
%  \autoref{sec:stlc} shows that we can achieve such reasoning easily 
%  by showing a type inference algorithm for \eqsat
%  with simply typed lambda calculus, 
%  something difficult to express in existing \eqsat frameworks.
% We go even beyond: the key idea of the type inference algorithm
%  for Hindley-Milner type systems is exactly its unification over types,
%  so there is nothing preventing us from doing it.
% \autoref{sec:hm} shows Hindley-Milner type inference in \egglog.
% Finally,
%  \autoref{sec:other-pearls} shows three more examples:
%  \egglog can solve multivariable equations by 
%  rewriting not only terms but also equations;
%  can encode, hashcons, and extract proofs for Datalog predicates;
%  and can use equational preconditions--guarded rewrite rules
%  to optimize matrix expressions with Kronecker products.

% \mw{point to supp material}
% \yz{we need an example to show analyses over eqsat}